\subsubsection*{1.3.4. 预备说明。}
设 $f:X\longrightarrow Y$ 是可微流形的光滑态射（按定义即无临界点的态射）。$X/Y$ 上的一个\underline{联络} $\nabla$ 是 $X$ 的切丛中与 $\ker(df)$ 互补的向量子丛。对 $Y$ 上的向量场 $u$，存在唯一的 $X$ 上向量场 $\nabla(u)$，提升 $u$ 并位于此子丛内。若
\[
\nabla([u,v])=[\nabla(u),\nabla(v)].
\]
则称此联络可积。当 $f$ 固有时，可积联络给出 $X/Y$ 在 $Y$ 上的局部平凡化
\[
\alpha:f^{-1}(U)\xrightarrow{\ \sim\ }X_0\times U,
\]
它与到 $U$ 的投影相容，并满足 $d\alpha(\nabla(u))=(0,u)$。当 $f$ 不固有时，若在无穷远处有适当条件，同一结论仍成立；例如，对某个适当的穷竭函数 $g$，$\nabla(u)g$ 恒为零即已足够。

\subsubsection*{1.3.5. Thom 同痕定理。}
\textit{设
\[
D=\{z\in\mathbb C:|z|<1\},
\qquad D^*=D-\{0\},
\]
且 $f:X\longrightarrow D$ 是解析空间的固有态射。将 $D$ 换为较小圆盘，并将 $X$ 换为其逆像后，$X^*=f^{-1}(D^*)$ 存在满足以下性质的分层 $\Sigma$：}

\begin{enumerate}[label=(\alph*),leftmargin=2.7em]
\item \textit{$\Sigma$ 是 $X^*$ 的一个由闭子集构成的滤过
\[
X^*=F_0\supset F_1\supset\cdots\supset F_n
\supset F_{n+1}=\varnothing
\]
使 $F_i$ 在 $X$ 中的闭包 $\overline F_i$ 是 $X$ 的解析子空间。}

\item \textit{开层片 $F_i-F_{i+1}$ 在 $D^*$ 上光滑。}

\item \textit{$\Sigma$ 满足 Whitney 条件 (a) 和 (b)。}
\end{enumerate}

\textit{对这样的分层 $\Sigma$，各 $F_i-F_{i+1}$ 上存在可积联络 $\nabla_i$，使得}

\begin{enumerate}[label=(\roman*),leftmargin=2.7em]
\item \textit{在 $D^*$ 上局部地，$\nabla_i$ 定义 $F_i-F_{i+1}$ 的平凡化；}

\item \textit{这些平凡化可粘合为 $X^*/D^*$ 的局部平凡化。}
\end{enumerate}

\textit{最后，若 $(Y_\alpha)$ 是 $X$ 的有限个闭解析子空间，则进一步缩小 $D$ 后，可选择 $\Sigma$ 使每个 $Y_\alpha$ 都是若干开层片 $(F_i-F_{i+1})$ 的连通分支的并。上述联络 $\nabla$ 于是局部平凡化 $(X,(Y_\alpha))$。}

\subsubsection*{1.3.6.}
当 $X$ 在 $D$ 上固有且满足下述条件时，同构 (1.3.3.2) 对以层复形 $K$ 为系数的超上同调仍成立：

\medskip
\noindent $(*)$ 对每个 $i\in\mathbb Z$，空间 $D^*$ 可由开子集 $U$ 覆盖，且存在与到 $U$ 的投影相容的同构
\[
\varphi:f^{-1}(U)\xrightarrow{\ \sim\ }Y\times U,
\]
以及 $Y$ 上的层 $G$，使在 $f^{-1}(U)$ 上，$\underline H^i(K)$ 同构于 $\varphi^*\operatorname{pr}_1^*G$。

在此假设下，1.2.6 的三角给出长正合列
\begin{equation}\tag{1.3.6.1}
\cdots\longrightarrow
\mathbb H^i(X_0,K)\longrightarrow
\mathbb H^i(X_t,K)\longrightarrow
\mathbb H^i(X_0,R\Phi_t(K))\longrightarrow\cdots.
\end{equation}
\noindent\emph{编校注。} 原印本此处引用三角“1.2.1”；所使用的区别三角是在 1.2.6 中构造的那个。

由 Thom 定理，只要 $X$ 具有复解析分层，且所有 $\underline H^i(K)$ 在每个开层片上局部常值，则（缩小 $D$ 后）条件 $(*)$ 成立。

